% GENERATED FILE. Edit canonical lessons, then run npm run generate:books.
% canonical-source-hash: fnv1a64:15e6ca25f6a5e770
% canonical-lessons: PA-C286-jhari, PA-C286-gajjna, PA-C286-varhna, PA-C286-jhakkhar, PA-C286-chamakna

\chapter{Rain on the Way}
\label{ch:pa-a2-rain-on-the-way}
\hlchaptermodality{286}

\begin{chapteropening}
\textbf{By the end of this chapter:} \emph{I can say what the rain and the storm will do.}

Complete the last lesson of chapter 286: I can say what the rain and the storm will do.
\end{chapteropening}

\section[\texorpdfstring{\pa{ਝੜੀ} (jhaṛī)}{ਝੜੀ (jhaṛī)}]{\pa{ਝੜੀ} (jhaṛī) — a spell of steady rain}

\label{lesson:PA-C286-jhari}



\begin{quote}
Before the new one: say the Punjabi for to get tired, then the Punjabi for dawn.
\end{quote}

\subsection*{You'll want to know: \pa{ਝੜੀ}}
\textbf{\pa{ਝੜੀ}} — \emph{jhaṛī} — \textquotedblleft{}a spell of steady rain\textquotedblright{}.

\begin{quote}
\pa{ਬਰਸਾਤ} \pa{ਵਿੱਚ} \pa{ਲੰਮੀ} \pa{ਝੜੀ} \pa{ਹੁੰਦੀ} \pa{ਹੈ।}
\end{quote}

\begin{quote}
\emph{barsāt vich lammī jhaṛī hundī hai.}
\end{quote}

\begin{quote}
In the rainy season there is a long spell of steady rain.
\end{quote}

A feminine noun.

\begin{cousinweb}
\textbf{\pa{ਝੜੀ}} (\emph{jhaṛī}) is inherited from a reconstructed \emph{*jhaḍī}, \textquotedblleft{}rainstorm\textquotedblright{}; CDIAL lists Punjabi \emph{jhaṛī}, \textquotedblleft{}continued rain\textquotedblright{}, under it (CDIAL 5329).
\end{cousinweb}

\begin{grammarlens}[title={What you've built}]
One more word for saying what comes next.
\end{grammarlens}

\subsection*{Your turn}
\begin{itemize}
\raggedright
  \item \emph{Say it:} \emph{jhaṛī}
  \item \emph{Say it:} \emph{jhaṛī}, once more
  \item \emph{Say it:} \emph{barsāt vich lammī jhaṛī hundī hai.} — In the rainy season there is a long spell of steady rain.
  \item \emph{Recall:} say the Punjabi for to get tired, then the Punjabi for dawn, then say \emph{jhaṛī} again
\end{itemize}

\begin{tcolorbox}[breakable,colback=teal!4,colframe=teal!35!black,title={Before you move on}]
What does \pa{ਝੜੀ} mean? (\textquotedblleft{}A spell of steady rain\textquotedblright{}.) Say it once more.
\end{tcolorbox}
\section[\texorpdfstring{\pa{ਗੱਜਣਾ} (gajjṇā)}{ਗੱਜਣਾ (gajjṇā)}]{\pa{ਗੱਜਣਾ} (gajjṇā) — to thunder; to roar}

\label{lesson:PA-C286-gajjna}



\begin{quote}
Before the new one: say the Punjabi for dawn, then the Punjabi for a spell of steady rain.
\end{quote}

\subsection*{You'll want to know: \pa{ਗੱਜਣਾ}}
\textbf{\pa{ਗੱਜਣਾ}} — \emph{gajjṇā} — \textquotedblleft{}to thunder; to roar\textquotedblright{}.

\begin{quote}
\pa{ਰਾਤ} \pa{ਨੂੰ} \pa{ਬੱਦਲ} \pa{ਗੱਜਣਗੇ।}
\end{quote}

\begin{quote}
\emph{rāt nū̃ baddal gajjaṇge.}
\end{quote}

\begin{quote}
The clouds will thunder at night.
\end{quote}

Past (it …): \pa{ਉਹ} \pa{ਬੱਦਲ} \pa{ਗੱਜਿਆ}, \emph{oh baddal gajjiā}.

\begin{cousinweb}
\textbf{\pa{ਗੱਜਣਾ}} (\emph{gajjṇā}) is inherited from Sanskrit \emph{garjati}, \textquotedblleft{}thunders, roars\textquotedblright{}; CDIAL lists Punjabi \emph{gajjṇā} under it (CDIAL 4046).
\end{cousinweb}

\begin{grammarlens}[title={What you've built}]
One more word for saying what comes next.
\end{grammarlens}

\subsection*{Your turn}
\begin{itemize}
\raggedright
  \item \emph{Say it:} \emph{gajjṇā}
  \item \emph{Say it:} \emph{gajjṇā}, once more
  \item \emph{Say it:} \emph{rāt nū̃ baddal gajjaṇge.} — The clouds will thunder at night.
  \item \emph{Recall:} say the Punjabi for dawn, then the Punjabi for a spell of steady rain, then say \emph{gajjṇā} again
  \item \emph{Write it:} copy \textbf{\pa{ਗੱਜਣਾ}} once, checking the whole word
\end{itemize}

\begin{tcolorbox}[breakable,colback=teal!4,colframe=teal!35!black,title={Before you move on}]
What does \pa{ਗੱਜਣਾ} mean? (\textquotedblleft{}To thunder; to roar\textquotedblright{}.) Say it once more.
\end{tcolorbox}
\section[\texorpdfstring{\pa{ਵਰ੍ਹਨਾ} (varhnā)}{ਵਰ੍ਹਨਾ (varhnā)}]{\pa{ਵਰ੍ਹਨਾ} (varhnā) — to pour down; to rain hard}

\label{lesson:PA-C286-varhna}



\begin{quote}
Before the new one: say the Punjabi for a spell of steady rain, then the Punjabi for to thunder.
\end{quote}

\subsection*{You'll want to know: \pa{ਵਰ੍ਹਨਾ}}
\textbf{\pa{ਵਰ੍ਹਨਾ}} — \emph{varhnā} — \textquotedblleft{}to pour down; to rain hard\textquotedblright{}.

\begin{quote}
\pa{ਕੱਲ੍ਹ} \pa{ਬਹੁਤ} \pa{ਮੀਂਹ} \pa{ਵਰ੍ਹੇਗਾ।}
\end{quote}

\begin{quote}
\emph{kall bahut mīṁh varhegā.}
\end{quote}

\begin{quote}
Tomorrow a lot of rain will pour down.
\end{quote}

Past (it …): \pa{ਉਹ} \pa{ਮੀਂਹ} \pa{ਵਰ੍ਹਿਆ}, \emph{oh mīṁh varhiā}.

\begin{cousinweb}
\textbf{\pa{ਵਰ੍ਹਨਾ}} (\emph{varhnā}) is inherited from Sanskrit \emph{varṣati}, \textquotedblleft{}it rains\textquotedblright{}; CDIAL lists Punjabi \emph{varhṇā}, \textquotedblleft{}to rain\textquotedblright{}, under it (CDIAL 11394).
\end{cousinweb}

\begin{grammarlens}[title={What you've built}]
One more word for saying what comes next.
\end{grammarlens}

\subsection*{Your turn}
\begin{itemize}
\raggedright
  \item \emph{Say it:} \emph{varhnā}
  \item \emph{Say it:} \emph{varhnā}, once more
  \item \emph{Say it:} \emph{kall bahut mīṁh varhegā.} — Tomorrow a lot of rain will pour down.
  \item \emph{Recall:} say the Punjabi for a spell of steady rain, then the Punjabi for to thunder, then say \emph{varhnā} again
  \item \emph{Write it:} copy \textbf{\pa{ਵਰ੍ਹਨਾ}} once, checking the whole word
\end{itemize}

\begin{tcolorbox}[breakable,colback=teal!4,colframe=teal!35!black,title={Before you move on}]
What does \pa{ਵਰ੍ਹਨਾ} mean? (\textquotedblleft{}To pour down; to rain hard\textquotedblright{}.) Say it once more.
\end{tcolorbox}
\section[\texorpdfstring{\pa{ਝੱਖੜ} (jhakkhaṛ)}{ਝੱਖੜ (jhakkhaṛ)}]{\pa{ਝੱਖੜ} (jhakkhaṛ) — a squall; a gale}

\label{lesson:PA-C286-jhakkhar}



\begin{quote}
Before the new one: say the Punjabi for to thunder, then the Punjabi for to pour down.
\end{quote}

\subsection*{You'll want to know: \pa{ਝੱਖੜ}}
\textbf{\pa{ਝੱਖੜ}} — \emph{jhakkhaṛ} — \textquotedblleft{}a squall; a gale\textquotedblright{}.

\begin{quote}
\pa{ਝੱਖੜ} \pa{ਨਾਲ} \pa{ਪੁਰਾਣਾ} \pa{ਰੁੱਖ} \pa{ਡਿੱਗੇਗਾ।}
\end{quote}

\begin{quote}
\emph{jhakkhaṛ nāl purāṇā rukkh ḍiggegā.}
\end{quote}

\begin{quote}
The old tree will fall in the squall.
\end{quote}

A masculine noun.

\begin{cousinweb}
\textbf{\pa{ਝੱਖੜ}} (\emph{jhakkhaṛ}) is inherited from a reconstructed \emph{*jhakk}, \textquotedblleft{}sudden movement or blast\textquotedblright{}; CDIAL lists Punjabi \emph{jhakkhaṛ}, \textquotedblleft{}squall\textquotedblright{}, under it (CDIAL 5316).
\end{cousinweb}

\begin{grammarlens}[title={What you've built}]
One more word for saying what comes next.
\end{grammarlens}

\subsection*{Your turn}
\begin{itemize}
\raggedright
  \item \emph{Say it:} \emph{jhakkhaṛ}
  \item \emph{Say it:} \emph{jhakkhaṛ}, once more
  \item \emph{Say it:} \emph{jhakkhaṛ nāl purāṇā rukkh ḍiggegā.} — The old tree will fall in the squall.
  \item \emph{Recall:} say the Punjabi for to thunder, then the Punjabi for to pour down, then say \emph{jhakkhaṛ} again
  \item \emph{Write it:} copy \textbf{\pa{ਝੱਖੜ}} once, checking the whole word
\end{itemize}

\begin{tcolorbox}[breakable,colback=teal!4,colframe=teal!35!black,title={Before you move on}]
What does \pa{ਝੱਖੜ} mean? (\textquotedblleft{}A squall; a gale\textquotedblright{}.) Say it once more.
\end{tcolorbox}
\section[\texorpdfstring{\pa{ਚਮਕਣਾ} (chamakṇā)}{ਚਮਕਣਾ (chamakṇā)}]{\pa{ਚਮਕਣਾ} (chamakṇā) — to flash; to shine}

\label{lesson:PA-C286-chamakna}



\begin{quote}
Before the new one: say the Punjabi for to pour down, then the Punjabi for a squall.
\end{quote}

\subsection*{You'll want to know: \pa{ਚਮਕਣਾ}}
\textbf{\pa{ਚਮਕਣਾ}} — \emph{chamakṇā} — \textquotedblleft{}to flash; to shine\textquotedblright{}.

\begin{quote}
\pa{ਬੱਦਲਾਂ} \pa{ਵਿੱਚ} \pa{ਬਿਜਲੀ} \pa{ਚਮਕੇਗੀ।}
\end{quote}

\begin{quote}
\emph{baddlāṁ vich bijlī chamkegī.}
\end{quote}

\begin{quote}
Lightning will flash in the clouds.
\end{quote}

Past (it …): \pa{ਉਹ} \pa{ਬਿਜਲੀ} \pa{ਚਮਕੀ}, \emph{oh bijlī chamkī}.

\begin{cousinweb}
\textbf{\pa{ਚਮਕਣਾ}} (\emph{chamakṇā}) is inherited from a reconstructed \emph{*camakka}, \textquotedblleft{}sudden movement\textquotedblright{}; CDIAL lists Punjabi \emph{camakṇā}, \textquotedblleft{}to shine\textquotedblright{}, under it (CDIAL 4676).
\end{cousinweb}

\begin{grammarlens}[title={What you've built}]
One more word for saying what comes next.
\end{grammarlens}

\subsection*{Your turn}
\begin{itemize}
\raggedright
  \item \emph{Say it:} \emph{chamakṇā}
  \item \emph{Say it:} \emph{chamakṇā}, once more
  \item \emph{Say it:} \emph{baddlāṁ vich bijlī chamkegī.} — Lightning will flash in the clouds.
  \item \emph{Recall:} say the Punjabi for to pour down, then the Punjabi for a squall, then say \emph{chamakṇā} again
  \item \emph{Write it:} copy \textbf{\pa{ਚਮਕਣਾ}} once, checking the whole word
\end{itemize}

\begin{tcolorbox}[breakable,colback=teal!4,colframe=teal!35!black,title={Before you move on}]
What does \pa{ਚਮਕਣਾ} mean? (\textquotedblleft{}To flash; to shine\textquotedblright{}.) Say it once more.
\end{tcolorbox}
